\section{Sim2Real y destilación de políticas}

\subsection{Estrategias de desplazamiento de distribución en Sim2Real}

Al migrar el entrenamiento de laboratorio al mundo físico, el Dynamics gap reduce significativamente el desempeño. Las contramedidas habituales incluyen: domain randomization (aleatorizar los parámetros del entorno), system identification (modelar la dinámica real) y una backup policy analítica como fallback. El `Sim-to-Real index` que la clase enfatiza se obtiene comparando la total variation distance entre el sim rollout y el real rollout, y da como resultado un indicador de policy drift.

\subsection{Destilación de políticas y despliegue de modelos pequeños}

El pipeline de distillation generalmente se divide en: teacher run -> dataset harvest -> student training -> student evaluation. Un dataset de calidad necesita cubrir los edge cases en los que el teacher se equivoca. En el student training suele usarse un objective compuesto de pérdida `mse` + `kl` para mantener la consistencia del comportamiento.

\begin{knowledgebox}{Por qué la distillation sigue siendo necesaria}
Aunque el teacher logre el desempeño más avanzado, su tamaño y su costo de inferencia elevado obligan a distillarlo hacia una policy latency-friendly. El proceso de destilación escribe el knowledge del teacher en el student en forma de soft target, de modo que el student pueda ejecutarse en un resource-constrained device.
\end{knowledgebox}

\subsection{Resumen de la sección}

Sim2Real requiere domain randomization, system ID y fallback policies. La destilación de políticas transfiere la capacidad del teacher a un student más ligero y es la vía clave para el despliegue en plataformas móviles/embebidas.

